% ============================================================
%  Guida completa: GMM con algoritmo EM
%  Corso: Graphical and Deep Learning (GDL) — Magistrale, Pisa
%  Midterm 1 — Fixed Structure + Incomplete Data
% ============================================================
\documentclass[11pt,a4paper]{article}

% ---- Packages ----
\usepackage[utf8]{inputenc}
\usepackage[T1]{fontenc}
\usepackage[italian]{babel}
\usepackage{amsmath}
\usepackage{amssymb}
\usepackage{booktabs}
\usepackage{enumitem}
\usepackage[margin=2.2cm]{geometry}
\usepackage{hyperref}
\usepackage{xcolor}
\usepackage[most]{tcolorbox}
\usepackage{algorithm}
\usepackage{algpseudocode}
\usepackage{listings}
\usepackage{array}

% ---- Hyperref setup ----
\hypersetup{
  colorlinks=true,
  linkcolor=blue!60!black,
  urlcolor=blue!60!black,
  citecolor=blue!60!black,
}

% ---- Listings setup for Python ----
\definecolor{codebg}{RGB}{245,245,245}
\definecolor{codegreen}{RGB}{0,128,0}
\definecolor{codepurple}{RGB}{128,0,128}
\definecolor{codeorange}{RGB}{204,102,0}

\lstdefinestyle{pythonstyle}{
  language=Python,
  backgroundcolor=\color{codebg},
  basicstyle=\ttfamily\small,
  keywordstyle=\color{blue}\bfseries,
  stringstyle=\color{codegreen},
  commentstyle=\color{gray}\itshape,
  numberstyle=\tiny\color{gray},
  numbers=left,
  numbersep=6pt,
  breaklines=true,
  frame=single,
  rulecolor=\color{gray!40},
  showstringspaces=false,
  tabsize=4,
  captionpos=b,
  morekeywords={self,np,True,False,None},
  emph={GaussianMixtureModel,__init__,fit,log_likelihood,bic,
        k_means_init,_compute_log_weighted},
  emphstyle=\color{codepurple}\bfseries,
  literate={à}{{\`a}}1 {è}{{\`e}}1 {ì}{{\`i}}1 {ò}{{\`o}}1 {ù}{{\`u}}1,
}
\lstset{style=pythonstyle}

% ---- tcolorbox environments ----
\newtcolorbox{infobox}[1][]{
  colback=blue!5!white,
  colframe=blue!60!black,
  fonttitle=\bfseries,
  title={#1},
  arc=2mm,
  boxrule=0.8pt,
}

\newtcolorbox{warnbox}[1][]{
  colback=red!5!white,
  colframe=red!70!black,
  fonttitle=\bfseries,
  title={#1},
  arc=2mm,
  boxrule=0.8pt,
}

\newtcolorbox{tipbox}[1][]{
  colback=green!5!white,
  colframe=green!50!black,
  fonttitle=\bfseries,
  title={#1},
  arc=2mm,
  boxrule=0.8pt,
}

\newtcolorbox{exbox}[1][]{
  colback=yellow!8!white,
  colframe=orange!70!black,
  fonttitle=\bfseries,
  title={#1},
  arc=2mm,
  boxrule=0.8pt,
}

% ---- Custom commands ----
\newcommand{\refslide}[2]{\textsc{#1}, slide~#2}

% ---- Title ----
\title{\textbf{Guida completa al codice:\\
Gaussian Mixture Model con algoritmo EM}\\[6pt]
\large Midterm 1 --- Graphical and Deep Learning}
\author{Appunti di studio}
\date{Anno Accademico 2025--2026}

\begin{document}
\maketitle
\tableofcontents
\newpage

% ============================================================
% SECTION 1: CONTESTO TEORICO
% ============================================================
\section{Contesto teorico}\label{sec:contesto}

Il corso di \emph{Graphical and Deep Learning} (GDL) organizza il problema
dell'apprendimento nei modelli grafici secondo due assi ortogonali
(\refslide{GDL7}{9}, \refslide{GDL8}{3}):

\begin{center}
\begin{tabular}{l|cc}
\toprule
 & \textbf{Complete Data} & \textbf{Incomplete Data} \\
\midrule
\textbf{Known Structure}   & MLE diretto          & \textbf{EM Algorithm} \\
\textbf{Unknown Structure}  & Structure Learning   & Structural EM \\
\bottomrule
\end{tabular}
\end{center}

\begin{infobox}[Posizione del midterm nella tassonomia]
Il nostro midterm si colloca nel quadrante \textbf{Fixed (Known) Structure +
Incomplete Data}. La struttura del modello grafico è fissata (un GMM con
$K$ componenti), ma i dati sono \emph{incompleti} perché non osserviamo le
variabili latenti~$z_j$ che indicano a quale componente appartiene ogni
campione~$x_j$. L'algoritmo per eccellenza in questo quadrante è
l'\textbf{Expectation-Maximization (EM)}.
\end{infobox}

\subsection{Perché i dati sono ``incompleti''?}

In un Gaussian Mixture Model, ogni campione $x_j$ è stato generato da
\emph{una} delle $K$ componenti gaussiane, ma non sappiamo quale. Se
conoscessimo l'assegnamento $z_j \in \{1,\ldots,K\}$ per ogni campione,
potremmo stimare i parametri con semplice MLE (quadrante in alto a sinistra).
Poiché $z_j$ è \emph{latente}, dobbiamo usare EM.

\subsection{Panoramica del flusso}

Il codice che abbiamo scritto segue esattamente questo schema:
\begin{enumerate}[label=\arabic*.]
  \item \textbf{Inizializzazione}: K-Means++ per le medie, varianza globale
        per le covarianze, pesi uniformi.
  \item \textbf{E-step}: calcolo delle responsabilità $\gamma_{jk}$
        (posterior sulle variabili latenti).
  \item \textbf{M-step}: aggiornamento di $\pi_k$, $\mu_k$, $\Sigma_k$
        usando le responsabilità.
  \item \textbf{Convergenza}: controllo sulla log-likelihood.
  \item \textbf{Model selection}: BIC per scegliere il numero ottimale
        di componenti.
\end{enumerate}

% ============================================================
% SECTION 2: IL MODELLO GMM
% ============================================================
\section{Il modello GMM}\label{sec:gmm}

\subsection{Definizione formale}

Un \emph{Gaussian Mixture Model} con $K$ componenti definisce la densità
di probabilità di un vettore $x \in \mathbb{R}^d$ come
(\refslide{GDL8}{19--20}):
\begin{equation}\label{eq:gmm}
  p(x \mid \theta) = \sum_{k=1}^{K} \pi_k \,
  \mathcal{N}(x \mid \mu_k, \Sigma_k)
\end{equation}
dove $\theta = \{\pi_k, \mu_k, \Sigma_k\}_{k=1}^{K}$ è l'insieme
dei parametri e:
\begin{itemize}[nosep]
  \item $\pi_k \ge 0$ sono i \textbf{mixing coefficients} (pesi),
        con $\sum_k \pi_k = 1$.
  \item $\mu_k \in \mathbb{R}^d$ è il \textbf{vettore medio} della
        $k$-esima componente.
  \item $\Sigma_k$ è la \textbf{matrice di covarianza} della $k$-esima
        componente.
\end{itemize}

\subsection{Gaussiana multivariata con covarianza diagonale}

La slide~18 di GDL8 mostra la gaussiana \textbf{univariata} ($d=1$,
un singolo numero reale):
\begin{equation}\label{eq:gauss_uni}
  \mathcal{N}(x \mid \mu, \sigma^2)
  = \frac{1}{\sqrt{2\pi}\;\sigma}
    \exp\!\left(-\frac{(x - \mu)^2}{2\sigma^2}\right)
\end{equation}

Questa formula ha tre pezzi: una costante $\frac{1}{\sqrt{2\pi}}$,
un fattore di scala $\frac{1}{\sigma}$ che dipende dalla ``larghezza''
della campana, e l'esponenziale che misura quanto $x$ \`{e} lontano
da~$\mu$.

\medskip

La \textbf{formula multivariata} ($d$ dimensioni) generalizza
ciascuno di questi tre pezzi:

\begin{center}
\renewcommand{\arraystretch}{1.6}
\begin{tabular}{c c c l}
\toprule
& \textbf{Univariata} ($d=1$) & \textbf{Multivariata} ($d$ dim.) &
\textbf{Perch\'{e}} \\
\midrule
Costante &
$\dfrac{1}{\sqrt{2\pi}}$ &
$\dfrac{1}{(2\pi)^{d/2}}$ &
una copia di $\sqrt{2\pi}$ per ogni dimensione \\[6pt]
Scala &
$\dfrac{1}{\sigma}$ &
$\dfrac{1}{|\Sigma|^{1/2}}$ &
$\sigma \to$ radice del determinante della matrice \\[6pt]
Distanza &
$\dfrac{(x-\mu)^2}{\sigma^2}$ &
$(x{-}\mu)^\top \Sigma^{-1}(x{-}\mu)$ &
distanza di Mahalanobis (tiene conto delle correlazioni) \\
\bottomrule
\end{tabular}
\end{center}

\medskip
Sostituendo nella stessa struttura della~\eqref{eq:gauss_uni} si
ottiene la formula generale:
\begin{equation}\label{eq:gauss_full}
  \mathcal{N}(x \mid \mu, \Sigma) = \frac{1}{(2\pi)^{d/2}
  \;|\Sigma|^{1/2}} \exp\!\left(-\frac{1}{2}(x-\mu)^\top
  \Sigma^{-1}(x-\mu)\right)
\end{equation}

Se poni $d=1$, $\Sigma = \sigma^2$ (scalare), $|\Sigma|^{1/2} = \sigma$
e $(x{-}\mu)^\top\Sigma^{-1}(x{-}\mu) = (x{-}\mu)^2/\sigma^2$:
si riottiene esattamente la~\eqref{eq:gauss_uni} della slide~18.

\begin{infobox}[Perch\'{e} $p(x_j \mid \mu_k, \Sigma_k) = \mathcal{N}(x_j \mid \mu_k, \Sigma_k)$?]
Questa uguaglianza \emph{non} \`{e} un risultato da dimostrare: \`{e} una
\textbf{scelta di modello}.  Quando definiamo un GMM, \emph{assumiamo} che
ogni componente~$k$ generi i dati secondo una distribuzione gaussiana.
Pertanto la densit\`{a} di probabilit\`{a} di osservare il campione~$x_j$,
sapendo che proviene dalla componente~$k$, \`{e} \emph{per definizione} la
funzione di densit\`{a} gaussiana:
\[
  p(x_j \mid z_j = k,\;\theta)
  \;=\;
  \mathcal{N}(x_j \mid \mu_k, \Sigma_k)
\]
In altre parole, $\mathcal{N}(\cdot)$ \`{e} semplicemente il \emph{nome}
che diamo alla funzione di densit\`{a} quando ha la forma a campana
definita dalla formula~\eqref{eq:gauss_full}.
Scrivere~$\mathcal{N}$ al posto di~$p$ \`{e} solo una notazione pi\`{u}
compatta: le due scritture indicano la stessa quantit\`{a}.
\end{infobox}

Nel nostro codice usiamo una \textbf{covarianza diagonale}:
$\Sigma_k = \mathrm{diag}(\sigma_{k,1}^2, \ldots, \sigma_{k,d}^2)$.
Questo semplifica enormemente i calcoli:
\begin{align}
  |\Sigma_k| &= \prod_{i=1}^{d} \sigma_{k,i}^2 \label{eq:det}\\
  (x-\mu_k)^\top \Sigma_k^{-1} (x-\mu_k) &=
  \sum_{i=1}^{d} \frac{(x_i - \mu_{k,i})^2}{\sigma_{k,i}^2}
  \label{eq:mahal}
\end{align}

\begin{tipbox}[Covarianza diagonale: vantaggi pratici]
Con covarianza diagonale:
\begin{itemize}[nosep]
  \item Non serve invertire matrici: basta dividere elemento per elemento.
  \item Il determinante è il prodotto delle varianze.
  \item La complessità scende da $O(d^3)$ a $O(d)$ per componente.
  \item In NumPy si memorizza come un vettore di shape
        \texttt{(n\_categories, n\_features)}, non come matrice $d\times d$.
\end{itemize}
\end{tipbox}

\subsection{Parametri nel codice e loro shape}

\begin{center}
\begin{tabular}{llll}
\toprule
\textbf{Parametro} & \textbf{Attributo} & \textbf{Shape NumPy} &
\textbf{Significato} \\
\midrule
$\pi_k$ & \texttt{self.\_pi} & \texttt{(K,)} &
Peso della componente $k$ \\
$\mu_k$ & \texttt{self.\_mu} & \texttt{(K, d)} &
Media della componente $k$ \\
$\sigma_k^2$ & \texttt{self.\_sigma} & \texttt{(K, d)} &
Varianze diagonali della componente $k$ \\
\bottomrule
\end{tabular}
\end{center}

Nel nostro caso: $K \in \{1,\ldots,6\}$ (provati tutti), $d=5$ features,
$N=800$ campioni dal file \texttt{train.csv}.

\begin{lstlisting}[caption={Costruttore della classe}]
class GaussianMixtureModel:
    def __init__(self, n_categories: int, n_features: int,
                 debug: bool = False, plot: bool = False):
        self.n_categories = n_categories
        self._pi = None    # Shape: (n_categories,)
        self._mu = None    # Shape: (n_categories, n_features)
        self._sigma = None # Shape: (n_categories, n_features)
        self.max_iteration = 1000
        self.epsilon = 1e-8
        self.debug = debug
        self.plot = plot
        self._ll_history = []
\end{lstlisting}

\begin{infobox}[Perche usare None e non array vuoti]
Inizializzare a \texttt{None} permette al metodo \texttt{fit()} di
controllare se i parametri sono già stati impostati dall'utente
(ad esempio per un warm-start) oppure se devono essere inizializzati
automaticamente. È un pattern comune nelle librerie ML come
\texttt{scikit-learn}.
\end{infobox}

% ============================================================
% SECTION 3: INCOMPLETE VS COMPLETE LIKELIHOOD
% ============================================================
\section{Incomplete vs Complete Likelihood}\label{sec:likelihood}

Questa distinzione è \textbf{fondamentale} per capire perché serve EM e
cosa calcola ogni pezzo del codice.

\subsection{Incomplete (Marginal) Likelihood}

La \textbf{incomplete likelihood} (\refslide{GDL8}{21}) è la probabilità
dei dati osservati, marginalizzando sulle variabili latenti:
\begin{equation}\label{eq:incomplete_ll}
  L(\theta \mid X) = \prod_{j=1}^{N} p(x_j \mid \theta)
  = \prod_{j=1}^{N} \sum_{k=1}^{K} \pi_k \,
  \mathcal{N}(x_j \mid \mu_k, \Sigma_k)
\end{equation}

In forma logaritmica (la \textbf{incomplete log-likelihood}):
\begin{equation}\label{eq:incomplete_logll}
  \ell(\theta \mid X) = \sum_{j=1}^{N} \log\!\left(
  \sum_{k=1}^{K} \pi_k \, \mathcal{N}(x_j \mid \mu_k, \Sigma_k)
  \right)
\end{equation}

\begin{warnbox}[Il logaritmo della somma non si semplifica]
Nella~\eqref{eq:incomplete_logll} il logaritmo ``avvolge'' una somma.
Non possiamo portare il $\log$ dentro la sommatoria. Per questo serve
il \textbf{log-sum-exp trick} (Sezione~\ref{sec:logsumexp}) e per
questo la MLE diretta è intrattabile --- si usa EM.
\end{warnbox}

\subsection{Complete Likelihood}

Se conoscessimo le variabili latenti $z_j$, avremmo la \textbf{complete
likelihood} (\refslide{GDL8}{24}). Si introduce un indicatore
$z_{jk} \in \{0,1\}$ con $\sum_k z_{jk}=1$ per ogni $j$:
\begin{equation}\label{eq:complete_ll}
  L_c(\theta \mid X, Z) = \prod_{j=1}^{N} \prod_{k=1}^{K}
  \left[\pi_k \, \mathcal{N}(x_j \mid \mu_k, \Sigma_k)\right]^{z_{jk}}
\end{equation}

La \textbf{complete log-likelihood}:
\begin{equation}\label{eq:complete_logll}
  \ell_c(\theta \mid X, Z) = \sum_{j=1}^{N} \sum_{k=1}^{K}
  z_{jk} \left[\log \pi_k + \log \mathcal{N}(x_j \mid \mu_k,
  \Sigma_k)\right]
\end{equation}

\begin{infobox}[Come si collegano nel codice]
\begin{itemize}[nosep]
  \item Il metodo \texttt{log\_likelihood()} calcola la
        \textbf{incomplete log-likelihood}~\eqref{eq:incomplete_logll}.
        È usato per il BIC e per il monitoraggio.
  \item L'E-step di EM ottimizza internamente la
        \textbf{expected complete log-likelihood}: sostituisce $z_{jk}$
        con il suo valore atteso $\gamma_{jk}$ (le responsabilità).
  \item L'M-step massimizza questa quantità rispetto a $\theta$.
\end{itemize}
\end{infobox}

\subsection{Il metodo log\_likelihood nel codice}

\begin{lstlisting}[caption={Calcolo della incomplete log-likelihood}]
def log_likelihood(self, samples: np.ndarray) -> np.ndarray:
    log_w = self._compute_log_weighted(samples)
    a_max = np.max(log_w, axis=1)
    log_sum = a_max + np.log(
        np.sum(np.exp(log_w - a_max[:, None]), axis=1)
    )
    return np.sum(log_sum)
\end{lstlisting}

Corrispondenza riga per riga con la formula~\eqref{eq:incomplete_logll}:

\begin{center}
\renewcommand{\arraystretch}{1.4}
\begin{tabular}{>{\ttfamily\small}p{6.5cm} p{7.5cm}}
\toprule
\textbf{Codice} & \textbf{Formula} \\
\midrule
log\_w = self.\_compute\_log\_weighted(samples) &
$\log\bigl(\pi_k \, \mathcal{N}(x_j \mid \mu_k, \Sigma_k)\bigr)$
per ogni $j,k$ \\
a\_max = np.max(log\_w, axis=1) &
$a_j = \max_k \log w_{jk}$ (per stabilità numerica) \\
np.exp(log\_w - a\_max[:, None]) &
$\exp(\log w_{jk} - a_j)$ \\
np.sum(..., axis=1) &
$\sum_k \exp(\log w_{jk} - a_j)$ \\
a\_max + np.log(...) &
$a_j + \log\sum_k \exp(\log w_{jk} - a_j)
= \log\sum_k w_{jk}$ \\
np.sum(log\_sum) &
$\sum_j \log\sum_k w_{jk} = \ell(\theta \mid X)$ \\
\bottomrule
\end{tabular}
\end{center}

Quindi \texttt{log\_likelihood} restituisce esattamente la
formula~\eqref{eq:incomplete_logll}, calcolata in modo numericamente
stabile grazie al log-sum-exp trick.

% ============================================================
% SECTION 4: _compute_log_weighted
% ============================================================
\section{\texttt{\_compute\_log\_weighted}: dalla formula al codice}
\label{sec:logweighted}

Questa è la funzione \emph{cuore} dell'implementazione. Calcola, per
ogni campione $j$ e ogni componente $k$:
\begin{equation}\label{eq:log_weighted}
  \log w_{jk} = \log \pi_k + \log \mathcal{N}(x_j \mid \mu_k, \Sigma_k)
\end{equation}

\subsection{Espansione della log-gaussiana}

Partendo dalla~\eqref{eq:gauss_full} con covarianza diagonale, il
logaritmo della gaussiana si decompone in tre termini:
\begin{equation}\label{eq:log_gauss_expanded}
  \log \mathcal{N}(x_j \mid \mu_k, \Sigma_k) =
  \underbrace{-\frac{d}{2}\log(2\pi)}_{\text{Termine A}}
  \underbrace{-\frac{1}{2}\sum_{i=1}^{d}\log\sigma_{k,i}^2}_{\text{Termine B}}
  \underbrace{-\frac{1}{2}\sum_{i=1}^{d}
  \frac{(x_{j,i}-\mu_{k,i})^2}{\sigma_{k,i}^2}}_{\text{Termine C}}
\end{equation}

\begin{infobox}[Derivazione passo-passo della formula~\eqref{eq:log_gauss_expanded}]

\textbf{Passo~0 --- riscrittura come prodotto.}
La formula~\eqref{eq:gauss_full} \`{e} gi\`{a} un prodotto di tre
pezzi; li separiamo solo per chiarezza:
\[
  \underbrace{
    \frac{1}{(2\pi)^{d/2} \; |\Sigma_k|^{1/2}}
  }_{\text{denominatore della frazione}}
  \;=\;
  \underbrace{\frac{1}{(2\pi)^{d/2}}}_{\text{fattore 1}}
  \;\cdot\;
  \underbrace{\frac{1}{|\Sigma_k|^{1/2}}}_{\text{fattore 2}}
\]
perch\'{e} $\;\dfrac{1}{A \cdot B} = \dfrac{1}{A}\cdot\dfrac{1}{B}$.
Quindi la~\eqref{eq:gauss_full} diventa:
\[
  \mathcal{N}(x_j \mid \mu_k,\Sigma_k)
  \;=\;
  \underbrace{\frac{1}{(2\pi)^{d/2}}}_{\text{fatt.\,1}}
  \;\cdot\;
  \underbrace{\frac{1}{|\Sigma_k|^{1/2}}}_{\text{fatt.\,2}}
  \;\cdot\;
  \underbrace{\exp\!\Bigl(-\tfrac{1}{2}(x_j-\mu_k)^\top
              \Sigma_k^{-1}(x_j-\mu_k)\Bigr)}_{\text{fatt.\,3}}
\]
Non abbiamo aggiunto n\'{e} tolto nulla: \`{e} la stessa espressione
scritta in forma pi\`{u} comoda per applicare il logaritmo.

\medskip
\textbf{Passo~1 --- applichiamo $\log$.}
Poich\'{e} $\log(a \cdot b \cdot c) = \log a + \log b + \log c$,
il logaritmo del prodotto diventa una \emph{somma} di tre logaritmi:

\begin{center}
\renewcommand{\arraystretch}{1.6}
\begin{tabular}{c p{5.8cm} c}
\toprule
\textbf{Fattore} & \textbf{Regola usata} & \textbf{Risultato} \\
\midrule
1 &
$\log\bigl((2\pi)^{-d/2}\bigr) = -\frac{d}{2}\log(2\pi)$
\newline usando $\log(a^n)=n\log a$
& Termine~A \\
2 &
$\log\bigl(|\Sigma_k|^{-1/2}\bigr) = -\frac{1}{2}\log|\Sigma_k|$
\newline $= -\frac{1}{2}\log\!\prod_i \sigma_{k,i}^2
         = -\frac{1}{2}\sum_i \log\sigma_{k,i}^2$
\newline usando $\log\!\prod = \sum\log$ (e $\Sigma_k$ diagonale,
\newline per cui $|\Sigma_k| = \prod_i\sigma_{k,i}^2$ come in~\eqref{eq:det})
& Termine~B \\
3 &
$\log\bigl(\exp(y)\bigr) = y$ \,($\log$ e $\exp$ si annullano);
\newline poi $(x_j{-}\mu_k)^\top\Sigma_k^{-1}(x_j{-}\mu_k)
       = \sum_i\frac{(x_{j,i}{-}\mu_{k,i})^2}{\sigma_{k,i}^2}$
\newline usando la~\eqref{eq:mahal} ($\Sigma_k$ diagonale)
& Termine~C \\
\bottomrule
\end{tabular}
\end{center}
\medskip
Sommando i tre risultati si ottiene esattamente la~\eqref{eq:log_gauss_expanded}.

\medskip
\textbf{Verifica rapida con $d=1$.}  Con una sola feature: Termine~A $= -\frac{1}{2}\log(2\pi)$,
Termine~B $= -\frac{1}{2}\log\sigma_k^2 = -\log\sigma_k$,
Termine~C $= -\frac{(x_j-\mu_k)^2}{2\sigma_k^2}$.
Sommando: $\log\!\left(\frac{1}{\sqrt{2\pi}\,\sigma_k}
\exp\!\bigl(-\frac{(x_j-\mu_k)^2}{2\sigma_k^2}\bigr)\right)$,
ossia il logaritmo della gaussiana univariata della slide~18.
\end{infobox}

\subsection{Il codice completo}

\begin{lstlisting}[caption={\_compute\_log\_weighted: il cuore del GMM}]
def _compute_log_weighted(self, samples: np.ndarray):
    d = samples.shape[1]
    differences = samples[:, None, :] - self._mu[None, :, :]
    differences_squares = differences ** 2
    scaled = differences_squares / self._sigma[None, :, :]
    third_term = np.sum(scaled, axis=2)
    second_term_log_var_sum = np.sum(np.log(self._sigma), axis=1)
    log_gauss = (-0.5 * d * np.log(2 * np.pi)
                 - 0.5 * second_term_log_var_sum[None, :]
                 - 0.5 * third_term)
    log_weighted = log_gauss + np.log(self._pi)[None, :]
    return log_weighted
\end{lstlisting}

\subsection{Analisi riga per riga}

\subsubsection{Riga 2: dimensionalità}
\begin{lstlisting}
d = samples.shape[1]  # d = 5 nel nostro caso
\end{lstlisting}
Semplicemente il numero di features. \texttt{samples} ha shape \texttt{(N, d)} = \texttt{(800, 5)}.

\subsubsection{Riga 3: differenze con broadcasting}
\begin{lstlisting}
differences = samples[:, None, :] - self._mu[None, :, :]
\end{lstlisting}

Questa è la riga più importante da capire. Il broadcasting di NumPy
calcola $x_{j,i} - \mu_{k,i}$ per \emph{tutte} le coppie $(j,k)$
simultaneamente:

\begin{center}
\renewcommand{\arraystretch}{1.3}
\begin{tabular}{lcl}
\toprule
\textbf{Operando} & \textbf{Shape} & \textbf{Significato degli assi} \\
\midrule
\texttt{samples[:, None, :]} & $(N, 1, d)$ &
asse 0 = campione $j$, asse 1 = broadcast, asse 2 = feature $i$ \\
\texttt{self.\_mu[None, :, :]} & $(1, K, d)$ &
asse 0 = broadcast, asse 1 = componente $k$, asse 2 = feature $i$ \\
\midrule
\texttt{differences} & $(N, K, d)$ &
$\text{differences}[j,k,i] = x_{j,i} - \mu_{k,i}$ \\
\bottomrule
\end{tabular}
\end{center}

\begin{exbox}[Esempio concreto di broadcasting]
Con $N=800$, $K=4$, $d=5$:
\begin{itemize}[nosep]
  \item \texttt{samples[:, None, :]} ha shape $(800, 1, 5)$ --- ogni
        campione viene ``ripetuto'' per $K$ componenti.
  \item \texttt{self.\_mu[None, :, :]} ha shape $(1, 4, 5)$ --- ogni
        media viene ``ripetuta'' per $N$ campioni.
  \item Il risultato ha shape $(800, 4, 5)$ --- tutte le $800 \times 4
        = 3200$ differenze, ciascuna un vettore di 5 elementi.
\end{itemize}
\end{exbox}

\subsubsection{Righe 4--5: distanza di Mahalanobis scalata}
\begin{lstlisting}
differences_squares = differences ** 2
scaled = differences_squares / self._sigma[None, :, :]
\end{lstlisting}

\begin{center}
\renewcommand{\arraystretch}{1.3}
\begin{tabular}{lccc}
\toprule
\textbf{Variabile} & \textbf{Shape input} & \textbf{Operazione} &
\textbf{Shape output} \\
\midrule
\texttt{differences\_squares} & $(N,K,d)$ & $(\cdot)^2$ & $(N,K,d)$ \\
\texttt{self.\_sigma[None,:,:]} & $(1,K,d)$ & broadcast & --- \\
\texttt{scaled} & $(N,K,d) \div (1,K,d)$ &
$\frac{(x_{j,i}-\mu_{k,i})^2}{\sigma_{k,i}^2}$ & $(N,K,d)$ \\
\bottomrule
\end{tabular}
\end{center}

\subsubsection{Riga 6: somma sulle features --- Termine C}
\begin{lstlisting}
third_term = np.sum(scaled, axis=2)  # Shape: (N, K)
\end{lstlisting}

Calcola:
\[
  \texttt{third\_term}[j,k] = \sum_{i=1}^{d}
  \frac{(x_{j,i}-\mu_{k,i})^2}{\sigma_{k,i}^2}
\]

Questa è la distanza di Mahalanobis al quadrato (con covarianza diagonale).
La somma su \texttt{axis=2} collassa la dimensione delle features,
passando da shape $(N,K,d)$ a $(N,K)$.

\subsubsection{Riga 7: log-determinante --- Termine B}
\begin{lstlisting}
second_term_log_var_sum = np.sum(np.log(self._sigma), axis=1)
\end{lstlisting}

Calcola per ogni componente $k$:
\[
  \texttt{second\_term\_log\_var\_sum}[k] = \sum_{i=1}^{d}
  \log \sigma_{k,i}^2 = \log \prod_{i=1}^{d} \sigma_{k,i}^2
  = \log |\Sigma_k|
\]

Shape: \texttt{(K,)}. Dalla~\eqref{eq:det}, il determinante di una
matrice diagonale è il prodotto degli elementi diagonali, quindi la
somma dei logaritmi è il logaritmo del determinante.

\subsubsection{Righe 8--10: assemblaggio della log-gaussiana}
\begin{lstlisting}
log_gauss = (-0.5 * d * np.log(2 * np.pi)         # Termine A
             - 0.5 * second_term_log_var_sum[None, :]  # Termine B
             - 0.5 * third_term)                       # Termine C
\end{lstlisting}

\begin{center}
\renewcommand{\arraystretch}{1.3}
\begin{tabular}{lccl}
\toprule
\textbf{Termine} & \textbf{Shape} & \textbf{Broadcasting} &
\textbf{Formula} \\
\midrule
A: \texttt{-0.5*d*np.log(2*np.pi)} & scalare & $\to (N,K)$ &
$-\frac{d}{2}\log(2\pi)$ \\
B: \texttt{-0.5*...[None,:]} & $(1,K)$ & $\to (N,K)$ &
$-\frac{1}{2}\log|\Sigma_k|$ \\
C: \texttt{-0.5*third\_term} & $(N,K)$ & --- &
$-\frac{1}{2}\sum_i\frac{(x_{j,i}-\mu_{k,i})^2}{\sigma_{k,i}^2}$ \\
\midrule
\texttt{log\_gauss} & $(N,K)$ & --- &
$\log\mathcal{N}(x_j \mid \mu_k,\Sigma_k)$ \\
\bottomrule
\end{tabular}
\end{center}

\subsubsection{Riga 11: aggiunta dei pesi}
\begin{lstlisting}
log_weighted = log_gauss + np.log(self._pi)[None, :]
\end{lstlisting}

\begin{center}
\renewcommand{\arraystretch}{1.3}
\begin{tabular}{lcl}
\toprule
\textbf{Operando} & \textbf{Shape} & \textbf{Formula} \\
\midrule
\texttt{log\_gauss} & $(N,K)$ &
$\log\mathcal{N}(x_j \mid \mu_k,\Sigma_k)$ \\
\texttt{np.log(self.\_pi)[None,:]} & $(1,K)$ &
$\log\pi_k$ (broadcast su $N$) \\
\midrule
\texttt{log\_weighted} & $(N,K)$ &
$\log\pi_k + \log\mathcal{N}(x_j \mid \mu_k,\Sigma_k)$ \\
\bottomrule
\end{tabular}
\end{center}

Il risultato finale è:
\[
  \texttt{log\_weighted}[j,k] = \log\!\bigl(\pi_k \,
  \mathcal{N}(x_j \mid \mu_k, \Sigma_k)\bigr)
\]

\subsection{Tabella riepilogativa delle shape}

\begin{center}
\renewcommand{\arraystretch}{1.3}
\small
\begin{tabular}{>{\ttfamily}l c l}
\toprule
\textbf{Variabile} & \textbf{Shape} & \textbf{Contenuto matematico} \\
\midrule
samples & $(N, d)$ & Dati di input $x_j$ \\
samples[:, None, :] & $(N, 1, d)$ & $x_j$ con asse per broadcast \\
self.\_mu[None, :, :] & $(1, K, d)$ & $\mu_k$ con asse per broadcast \\
differences & $(N, K, d)$ & $x_{j,i} - \mu_{k,i}$ \\
differences\_squares & $(N, K, d)$ & $(x_{j,i} - \mu_{k,i})^2$ \\
self.\_sigma[None, :, :] & $(1, K, d)$ & $\sigma_{k,i}^2$ per broadcast \\
scaled & $(N, K, d)$ & $(x_{j,i}-\mu_{k,i})^2 / \sigma_{k,i}^2$ \\
third\_term & $(N, K)$ & $\sum_i (x_{j,i}-\mu_{k,i})^2/\sigma_{k,i}^2$ \\
second\_term\_log\_var\_sum & $(K,)$ & $\sum_i \log\sigma_{k,i}^2$ \\
log\_gauss & $(N, K)$ & $\log\mathcal{N}(x_j \mid \mu_k, \Sigma_k)$ \\
np.log(self.\_pi) & $(K,)$ & $\log\pi_k$ \\
log\_weighted & $(N, K)$ & $\log(\pi_k \cdot \mathcal{N}(x_j|\mu_k,\Sigma_k))$ \\
\bottomrule
\end{tabular}
\end{center}

% ============================================================
% SECTION 5: LOG-SUM-EXP TRICK
% ============================================================
\section{Log-Sum-Exp Trick}\label{sec:logsumexp}

\subsection{Il problema numerico}

Dobbiamo calcolare:
\[
  \log \sum_{k=1}^{K} w_{jk}
  \quad\text{dove}\quad
  w_{jk} = \pi_k \, \mathcal{N}(x_j \mid \mu_k, \Sigma_k)
\]

Il calcolo diretto richiede di esponenziare i
\texttt{log\_weighted}$_{jk}$ per ottenere $w_{jk}$, sommarli,
e poi prendere il logaritmo. Ma $\log w_{jk}$ può essere molto
negativo (es.\ $-500$), rendendo $w_{jk} = e^{-500} \approx 0$
in floating point (underflow).

\subsection{La soluzione matematica}

Il trucco si basa sull'identità:
\begin{equation}\label{eq:logsumexp}
  \log \sum_k e^{a_k} = a^* + \log \sum_k e^{a_k - a^*}
  \qquad\text{dove}\quad a^* = \max_k a_k
\end{equation}

\textbf{Dimostrazione}:
\[
  \log \sum_k e^{a_k} = \log\!\left(e^{a^*} \sum_k
  e^{a_k - a^*}\right) = a^* + \log \sum_k e^{a_k - a^*}
\]

\begin{tipbox}[Perche funziona]
Dopo la sottrazione di $a^*$, tutti gli esponenti $a_k - a^*$ sono
$\le 0$. Il termine con $a_k = a^*$ contribuisce $e^0 = 1$, quindi
la somma è $\ge 1$ e il logaritmo è $\ge 0$. Nessun underflow!
\end{tipbox}

\subsection{Esempio numerico}

Supponiamo $K=3$ e per un certo campione $j$:
\[
  \log w_{j,1} = -502, \quad
  \log w_{j,2} = -500, \quad
  \log w_{j,3} = -505
\]

\textbf{Calcolo diretto} (FALLISCE):
\[
  e^{-502} + e^{-500} + e^{-505} \approx 0 + 0 + 0 = 0
  \quad\Rightarrow\quad \log(0) = -\infty \quad\text{(errore!)}
\]

\textbf{Con log-sum-exp} (FUNZIONA):
\begin{align*}
  a^* &= \max(-502, -500, -505) = -500 \\
  &-500 + \log\!\left(e^{-502-(-500)} + e^{-500-(-500)}
  + e^{-505-(-500)}\right) \\
  &= -500 + \log\!\left(e^{-2} + e^{0} + e^{-5}\right) \\
  &= -500 + \log\!\left(0.1353 + 1 + 0.0067\right) \\
  &= -500 + \log(1.1420) \\
  &= -500 + 0.1328 = -499.867
\end{align*}

\subsection{Corrispondenza nel codice}

\begin{lstlisting}[caption={Log-sum-exp in log\_likelihood}]
log_w = self._compute_log_weighted(samples)  # (N, K)
a_max = np.max(log_w, axis=1)                # (N,)
log_sum = a_max + np.log(
    np.sum(np.exp(log_w - a_max[:, None]), axis=1)
)                                             # (N,)
return np.sum(log_sum)                        # scalare
\end{lstlisting}

\begin{center}
\renewcommand{\arraystretch}{1.3}
\begin{tabular}{>{\ttfamily\small}l l c}
\toprule
\textbf{Codice} & \textbf{Formula} & \textbf{Shape} \\
\midrule
log\_w & $a_{jk} = \log w_{jk}$ & $(N,K)$ \\
a\_max & $a_j^* = \max_k a_{jk}$ & $(N,)$ \\
log\_w - a\_max[:, None] & $a_{jk} - a_j^*$ & $(N,K)$ \\
np.exp(...) & $e^{a_{jk} - a_j^*}$ & $(N,K)$ \\
np.sum(..., axis=1) & $\sum_k e^{a_{jk} - a_j^*}$ & $(N,)$ \\
a\_max + np.log(...) & $a_j^* + \log\sum_k e^{a_{jk}-a_j^*}$ & $(N,)$ \\
np.sum(log\_sum) & $\sum_j \log\sum_k w_{jk}$ & scalare \\
\bottomrule
\end{tabular}
\end{center}

\begin{warnbox}[Lo stesso pattern si ripete nel fit]
Il log-sum-exp trick appare \emph{due volte} nel codice:
\begin{enumerate}[nosep]
  \item In \texttt{log\_likelihood()}: per calcolare $\ell(\theta|X)$.
  \item Nell'E-step di \texttt{fit()}: per calcolare il denominatore
        delle responsabilità $\gamma_{jk}$ (stesse righe identiche).
\end{enumerate}
\end{warnbox}

% ============================================================
% SECTION 6: K-MEANS++ INITIALIZATION
% ============================================================
\section{K-Means++ Initialization}\label{sec:kmeans}

\subsection{Perché è necessaria una buona inizializzazione}

L'algoritmo EM è garantito convergere a un \textbf{massimo locale}
della likelihood (\refslide{GDL8}{15}), ma non necessariamente al
massimo globale. Il grafico a slide~15 mostra come EM cresca
monotonicamente ma possa convergere a soluzioni diverse a seconda
del punto di partenza.

Una cattiva inizializzazione può portare a:
\begin{itemize}[nosep]
  \item Componenti ``vuote'' (nessun campione assegnato).
  \item Componenti sovrapposte (ridondanti).
  \item Convergenza lenta.
  \item Log-likelihood finale bassa (massimo locale scarso).
\end{itemize}

\subsection{L'algoritmo K-Means++}

K-Means++ sceglie i centroidi iniziali in modo che siano ben
distanziati tra loro:

\begin{algorithm}[H]
\caption{K-Means++ Initialization}
\begin{algorithmic}[1]
\State Scegli $\mu_1$ uniformemente a caso tra i campioni
\For{$k = 2, \ldots, K$}
  \State Per ogni campione $x_j$, calcola $D(x_j) = \min_{k' < k}
         \|x_j - \mu_{k'}\|^2$
  \State Scegli $\mu_k = x_j$ con probabilità proporzionale a $D(x_j)$
\EndFor
\end{algorithmic}
\end{algorithm}

\subsection{Il codice}

\begin{lstlisting}[caption={K-Means++ initialization}]
def k_means_init(self, samples: np.ndarray):
    sample_quantity, features_quantity = samples.shape
    centroids = np.empty((self.n_categories, features_quantity))
    centroids[0] = samples[(np.random.randint(sample_quantity))]
    for k in range(1, self.n_categories):
        distances_sq = np.array(
            [np.sum((samples - mu_k_prime)**2, axis=1)
             for mu_k_prime in centroids[:k]]
        )
        min_dist_sq = np.min(distances_sq, axis=0)
        probabilities = min_dist_sq / np.sum(min_dist_sq)
        next_index = np.random.choice(sample_quantity,
                                       p=probabilities)
        centroids[k] = samples[next_index]
    return centroids
\end{lstlisting}

Analisi:
\begin{itemize}
  \item \textbf{Riga 4}: il primo centroide è un campione scelto
        uniformemente a caso.
  \item \textbf{Righe 6--8}: per ogni centroide già scelto $\mu_{k'}$,
        calcola la distanza euclidea al quadrato da ogni campione.
        \texttt{distances\_sq} ha shape $(k, N)$.
  \item \textbf{Riga 9}: per ogni campione, prende la distanza minima
        dal centroide più vicino. Shape: $(N,)$.
  \item \textbf{Riga 10}: normalizza per ottenere una distribuzione
        di probabilità. I campioni più lontani dai centroidi esistenti
        hanno probabilità maggiore di essere scelti.
  \item \textbf{Righe 11--12}: campiona il prossimo centroide secondo
        questa distribuzione.
\end{itemize}

\begin{tipbox}[Intuizione geometrica]
K-Means++ garantisce che i centroidi iniziali siano ``sparsi'' nello
spazio. Il primo è casuale; il secondo sarà probabilmente lontano dal
primo; il terzo sarà lontano dai primi due; e così via. Questo evita
il problema di inizializzare due componenti nello stesso cluster.
\end{tipbox}

% ============================================================
% SECTION 7: INITIALIZATION IN FIT
% ============================================================
\section{Inizializzazione di pi, mu, sigma in fit}\label{sec:init}

\begin{lstlisting}[caption={Inizializzazione dei parametri nel metodo fit}]
def fit(self, samples: np.ndarray):
    if self._mu is None:
        self._mu = self.k_means_init(samples)
    if self._pi is None:
        self._pi = np.full(self.n_categories,
                           1.0 / self.n_categories)
    if self._sigma is None:
        self._sigma = np.tile(np.var(samples, axis=0),
                              (self.n_categories, 1))
\end{lstlisting}

\subsection{Medie: K-Means++}
Se \texttt{self.\_mu} non è stato pre-impostato, si usa K-Means++
(Sezione~\ref{sec:kmeans}). Output: shape $(K, d)$.

\subsection{Pesi: uniformi}
\begin{lstlisting}
self._pi = np.full(self.n_categories, 1.0 / self.n_categories)
\end{lstlisting}

\texttt{np.full(K, 1/K)} crea un vettore di $K$ elementi tutti uguali
a $1/K$. Per $K=4$: $\pi = [0.25, 0.25, 0.25, 0.25]$.

Questa è la scelta più ``agnostica'': nessuna componente è favorita
a priori. I pesi verranno aggiornati dall'M-step.

\subsection{Varianze: varianza globale replicata}
\begin{lstlisting}
self._sigma = np.tile(np.var(samples, axis=0),
                      (self.n_categories, 1))
\end{lstlisting}

Passo per passo:
\begin{enumerate}[nosep]
  \item \texttt{np.var(samples, axis=0)}: calcola la varianza di ogni
        feature su tutti i campioni. Shape: $(d,)$.
  \item \texttt{np.tile(..., (K, 1))}: replica il vettore $K$ volte,
        creando una matrice $(K, d)$ dove ogni riga è identica.
\end{enumerate}

\begin{infobox}[Perche la varianza globale e non casuale]
Inizializzare con la varianza globale dei dati è una scelta ragionevole:
\begin{itemize}[nosep]
  \item È dell'ordine di grandezza giusto (non troppo grande, non
        troppo piccola).
  \item Ogni componente parte con la stessa ``apertura'', e l'EM
        le specializza.
  \item Inizializzare con varianze casuali potrebbe creare componenti
        troppo strette (che catturano un solo punto) o troppo larghe
        (che ignorano la struttura).
\end{itemize}
\end{infobox}

% ============================================================
% SECTION 8: EM ALGORITHM
% ============================================================
\section{Algoritmo EM}\label{sec:em}

L'Expectation-Maximization \`{e} un algoritmo iterativo per la stima ML
con dati incompleti (\refslide{GDL8}{14}).

\subsection{Schema generale (dalle slide)}

L'EM alterna due passi (\refslide{GDL8}{14}):
\begin{enumerate}
  \item \textbf{E-step}: calcola $Q(\theta \mid \theta^{(t)}) =
        \mathbb{E}_{Z|X,\theta^{(t)}}[\ell_c(\theta \mid X, Z)]$
  \item \textbf{M-step}: $\theta^{(t+1)} = \arg\max_\theta
        Q(\theta \mid \theta^{(t)})$
\end{enumerate}

\begin{infobox}[La funzione $Q$ delle slide e il legame con il codice]

Nelle slide la funzione ausiliaria $Q$ \`{e} scritta cos\`{i}
(\refslide{GDL8}{25}):
\begin{equation}\label{eq:Q}
  Q(\theta \mid \theta^{(t)})
  = \sum_{j=1}^{N}\sum_{m=1}^{M}
    \underbrace{P(z_j = m \mid x_j,\;\theta^{(t)})}_{\text{peso}}
    \;\cdot\;
    \underbrace{\log\!\bigl[
      \pi_m^{(t)}\;
      P(x_j \mid \mu_m^{(t)}, \sigma_m^{(t)})
    \bigr]}_{\text{log del termine pesato}}
\end{equation}

Questa formula contiene \textbf{due ingredienti} distinti.
Vediamoli uno alla volta e colleghiamoli al codice.

\bigskip
\textbf{Ingrediente 1: il ``log del termine pesato''}
(\texttt{log\_w} nel codice).

Dentro il $\log[\cdots]$ della~\eqref{eq:Q} c'\`{e}:
\[
  \pi_m \cdot P(x_j \mid \mu_m, \sigma_m)
\]
Poich\'{e} nel GMM ogni componente genera dati gaussiani,
$P(x_j \mid \mu_m, \sigma_m) = \mathcal{N}(x_j \mid \mu_m, \Sigma_m)$
(come spiegato nella Sezione~2.2). Quindi:
\[
  \log\!\bigl[\pi_m \cdot P(x_j|\mu_m,\sigma_m)\bigr]
  = \log\pi_m + \log\mathcal{N}(x_j \mid \mu_m, \Sigma_m)
\]
Questo \`{e} esattamente ci\`{o} che calcola
\texttt{\_compute\_log\_weighted} (Sezione~\ref{sec:logweighted}):

\medskip
\begin{center}
\fbox{
$\texttt{log\_w}[j,\,m]
 \;=\; \log\pi_m + \log\mathcal{N}(x_j \mid \mu_m, \Sigma_m)
 \;=\; \text{il termine dentro il } \log[\cdots]
       \text{ della~\eqref{eq:Q}}$
}
\end{center}

\bigskip
\textbf{Ingrediente 2: il ``peso'' $P(z_j = m \mid x_j, \theta^{(t)})$}
(\texttt{gamma} nel codice).

Il primo fattore della doppia sommatoria \`{e} il
\textbf{posterior}: ``con i parametri attuali $\theta^{(t)}$,
qual \`{e} la probabilit\`{a} che il campione~$x_j$ sia stato generato
dalla componente~$m$?''

Questo posterior ha un nome pi\`{u} corto usato universalmente nei
libri di testo: si chiama \textbf{responsabilit\`{a}} e si indica
con $\gamma_{jm}$:
\[
  \gamma_{jm}
  \;\stackrel{\text{def}}{=}\;
  P(z_j = m \mid x_j,\;\theta^{(t)})
\]
Non \`{e} una formula diversa --- \`{e} solo un'\textbf{abbreviazione}
per non riscrivere ogni volta $P(z_j = m \mid x_j, \theta^{(t)})$.

\bigskip
\textbf{In sintesi, la $Q$ delle slide diventa:}
\[
  Q(\theta \mid \theta^{(t)})
  = \sum_{j=1}^{N}\sum_{m=1}^{M}
    \underbrace{\gamma_{jm}}_{\texttt{gamma[j,m]}}
    \;\cdot\;
    \underbrace{\log\!\bigl[\pi_m\,\mathcal{N}(x_j|\mu_m,\Sigma_m)\bigr]}
              _{\texttt{log\_w[j,m]}}
\]

\medskip
\begin{center}
\renewcommand{\arraystretch}{1.4}
\begin{tabular}{lll}
\toprule
\textbf{Slide} & \textbf{Simbolo breve} & \textbf{Codice} \\
\midrule
$P(z_j{=}m \mid x_j, \theta^{(t)})$
  & $\gamma_{jm}$ & \texttt{gamma[j, m]} \\
$\log\bigl[\pi_m\,P(x_j|\mu_m,\sigma_m)\bigr]$
  & $\log w_{jm}$ & \texttt{log\_w[j, m]} \\
\bottomrule
\end{tabular}
\end{center}
\end{infobox}

Per il GMM, l'E-step si riduce al calcolo delle \textbf{responsabilit\`{a}}
$\gamma_{jm}$, e l'M-step alle formule chiuse per aggiornare i parametri.

\subsection{E-step: calcolo delle responsabilit\`{a}}

\begin{infobox}[Come si calcola $\gamma_{jm}$ a partire dalla formula delle slide]
Nelle slide, $P(z_j = m \mid x_j, \theta^{(t)})$ \`{e} un posterior.
Per calcolarlo si usa il \textbf{teorema di Bayes}:
\[
  P(z_j{=}m \mid x_j)
  = \frac{P(x_j \mid z_j{=}m)\;P(z_j{=}m)}
         {P(x_j)}
  = \frac{P(x_j \mid z_j{=}m)\;\pi_m}
         {\sum_{m'=1}^{M} P(x_j \mid z_j{=}m')\;\pi_{m'}}
\]
Ora, $P(x_j \mid z_j{=}m) = \mathcal{N}(x_j \mid \mu_m, \Sigma_m)$
(scelta di modello), e $P(z_j{=}m) = \pi_m$ (il mixing coefficient).
Sostituendo:
\begin{equation}\label{eq:gamma}
  \gamma_{jm}
  = \frac{\pi_m \;\mathcal{N}(x_j \mid \mu_m, \Sigma_m)}
         {\sum_{m'=1}^{M} \pi_{m'}\;\mathcal{N}(x_j \mid \mu_{m'}, \Sigma_{m'})}
\end{equation}

Osserva che:
\begin{itemize}[nosep]
  \item Il \textbf{numeratore} \`{e} $\pi_m \cdot \mathcal{N}(x_j|\mu_m,\Sigma_m)$,
        ossia \emph{proprio} $w_{jm} = \exp(\texttt{log\_w}[j,m])$.
  \item Il \textbf{denominatore} \`{e} la somma $\sum_{m'} w_{jm'}$,
        cio\`{e} la somma del numeratore su tutte le componenti.
\end{itemize}

Quindi $\gamma_{jm}$ \`{e} semplicemente il numeratore diviso per la
somma dei numeratori:
\[
  \gamma_{jm} = \frac{w_{jm}}{\sum_{m'} w_{jm'}}
\]
\end{infobox}

\begin{lstlisting}[caption={E-step nel codice}]
# E-step
log_w = self._compute_log_weighted(samples)        # (N, K)
a_max = np.max(log_w, axis=1)                      # (N,)
log_denominator = a_max + np.log(
    np.sum(np.exp(log_w - a_max[:, None]), axis=1)
)                                                   # (N,)
gamma = np.exp(log_w - log_denominator[:, None])    # (N, K)
\end{lstlisting}

Corrispondenza con la formula~\eqref{eq:gamma}:

\begin{center}
\renewcommand{\arraystretch}{1.4}
\begin{tabular}{>{\ttfamily\small}p{6.5cm} p{7.5cm}}
\toprule
\textbf{Codice} & \textbf{Formula delle slide} \\
\midrule
log\_w[j,m] &
$\log\bigl(\pi_m \cdot P(x_j|\mu_m,\sigma_m)\bigr)$
= logaritmo del numeratore della~\eqref{eq:gamma} \\
log\_denominator[j] &
$\log\sum_{m'} \pi_{m'} P(x_j|\mu_{m'},\sigma_{m'})$
= logaritmo del denominatore (via log-sum-exp) \\
log\_w - log\_denominator[:, None] &
$\log\gamma_{jm} = \log(\text{num}) - \log(\text{den})$
(perch\'{e} $\log\frac{a}{b} = \log a - \log b$) \\
gamma = np.exp(...) &
$\gamma_{jm} = P(z_j{=}m \mid x_j, \theta^{(t)})$ \\
\bottomrule
\end{tabular}
\end{center}

\begin{tipbox}[Perch\'{e} si lavora in log-scala]
Non calcoliamo $\gamma$ con una divisione diretta
($\frac{\text{num}}{\text{den}}$) perch\'{e} i valori
$\pi_m \cdot \mathcal{N}(\ldots)$ possono essere numeri
microscopici (es.\ $10^{-300}$) che il computer arrotonda a~0
(underflow). Invece:
\begin{enumerate}[nosep]
  \item calcoliamo tutto in \emph{log-scala} (dove $-300$ \`{e}
        un numero perfettamente rappresentabile);
  \item usiamo il \textbf{log-sum-exp trick} (Sezione~\ref{sec:logsumexp})
        per il denominatore;
  \item alla fine facciamo $\exp(\log\text{num} - \log\text{den})$
        per tornare alle probabilit\`{a}.
\end{enumerate}
Le $\gamma_{jm}$ risultanti sono in $[0,1]$ e ogni riga somma a~1.
\end{tipbox}

\begin{infobox}[Il codice non calcola $Q$ esplicitamente --- e va bene cos\`{i}]
Potresti chiederti: ``ma nel codice dove viene calcolata la doppia
sommatoria della $Q$?'' La risposta \`{e}: \textbf{da nessuna parte},
e non serve.

L'unico motivo per cui la~$Q$ esiste nella teoria \`{e} giustificare
\emph{perch\'{e}} le formule dell'M-step funzionano.
In pratica basta calcolare $\gamma_{jm}$ (E-step) e poi usarlo
nelle formule chiuse dell'M-step (Sezione~\ref{sec:em_mstep}).
La $Q$ non viene mai valutata numericamente.
\end{infobox}

\subsection{M-step: aggiornamento dei parametri}\label{sec:em_mstep}

Le formule dell'M-step (\refslide{GDL8}{27}) sono:

\begin{align}
  N_k &= \sum_{j=1}^{N} \gamma_{jk} \label{eq:Nk}\\
  \pi_k^{\text{new}} &= \frac{N_k}{N} \label{eq:pi_update}\\
  \mu_k^{\text{new}} &= \frac{1}{N_k}\sum_{j=1}^{N}
  \gamma_{jk}\, x_j \label{eq:mu_update}\\
  \Sigma_k^{\text{new}} &= \frac{1}{N_k}\sum_{j=1}^{N}
  \gamma_{jk}\,(x_j - \mu_k^{\text{new}})
  (x_j - \mu_k^{\text{new}})^\top \label{eq:sigma_update}
\end{align}

Nel caso diagonale, la~\eqref{eq:sigma_update} diventa:
\begin{equation}\label{eq:sigma_diag}
  \sigma_{k,i}^{2,\text{new}} = \frac{1}{N_k}\sum_{j=1}^{N}
  \gamma_{jk}\,(x_{j,i} - \mu_{k,i}^{\text{new}})^2
\end{equation}

\begin{infobox}[Dove \`{e} finita la radice quadrata della slide?]
La slide~27 di GDL8 potrebbe mostrare la formula di aggiornamento
per la \textbf{deviazione standard} $\sigma_m$:
\[
  \sigma_m^{\text{new}}
  = \sqrt{\frac{1}{N_m}\sum_{j=1}^{N}
    \gamma_{jm}\,(x_j - \mu_m^{\text{new}})^2\,}
\]
La formula~\eqref{eq:sigma_diag} nella guida \`{e} per la
\textbf{varianza} $\sigma_{k,i}^{2}$, cio\`{e} il quadrato della
deviazione standard. La relazione \`{e} banale:
\[
  \sigma_m^2 = \left(\sigma_m\right)^2
  \qquad\Longleftrightarrow\qquad
  \sigma_m = \sqrt{\sigma_m^2}
\]
Se prendi la~\eqref{eq:sigma_diag} e applichi la radice quadrata
a entrambi i lati, riottieni esattamente la formula della slide.
\textbf{Sono la stessa formula} --- la slide mostra $\sigma$, la
guida mostra $\sigma^2$.

\medskip
\textbf{Perch\'{e} nel codice memorizziamo $\sigma^2$ e non
$\sigma$?}
Perch\'{e} ovunque nel modello serve la \emph{varianza}, mai la
deviazione standard:
\begin{itemize}[nosep]
  \item Nella gaussiana~\eqref{eq:gauss_full}: il denominatore
        contiene $|\Sigma|^{1/2} = \prod_i \sigma_{k,i}^2$
        sotto radice, e l'esponente contiene
        $\frac{(x_i-\mu_i)^2}{\sigma_{k,i}^2}$.
  \item Nel codice (\texttt{\_compute\_log\_weighted}):
        \texttt{self.\_sigma} appare in divisioni (\texttt{diff**2
        / self.\_sigma}) e in logaritmi (\texttt{np.log(self.\_sigma)})
        --- sempre come varianza.
\end{itemize}
Prendere una radice quadrata per poi riquadrarla ogni volta sarebbe
inutile. Quindi il codice salta la radice e lavora direttamente con
$\sigma^2$.
\end{infobox}

\begin{lstlisting}[caption={M-step nel codice}]
# M-step
N_k = gamma.sum(axis=0)             # (K,)
N_k = np.maximum(N_k, 1e-10)        # floor per sicurezza
self._pi = N_k / N                  # (K,)
self._mu = gamma.T @ samples / N_k[:, None]  # (K, d)
diff = samples[:, None, :] - self._mu[None, :, :]  # (N, K, d)
self._sigma = (np.sum(gamma[:, :, None] * diff**2, axis=0)
               / N_k[:, None])       # (K, d)
self._sigma = np.maximum(self._sigma, 1e-6)
\end{lstlisting}

\subsubsection{Riga 2: conteggio soft $N_k$}
\begin{lstlisting}
N_k = gamma.sum(axis=0)  # (K,)
\end{lstlisting}
Implementa la~\eqref{eq:Nk}. $N_k$ è il ``numero effettivo'' di
campioni assegnati alla componente $k$. Non è un intero: è una somma
di probabilità. $\sum_k N_k = N$.

\subsubsection{Riga 3: floor su $N_k$}
\begin{lstlisting}
N_k = np.maximum(N_k, 1e-10)
\end{lstlisting}

\begin{warnbox}[Perche il floor su N\_k]
Se una componente non ha responsabilità su nessun campione
($N_k \approx 0$), la divisione per $N_k$ nelle formule successive
causerebbe divisione per zero o valori astronomici. Il floor a
$10^{-10}$ previene questo problema. Nella pratica, se $N_k$ è
così piccolo, la componente è ``morta'' e i suoi parametri diventano
irrilevanti.
\end{warnbox}

\subsubsection{Riga 4: aggiornamento pesi $\pi_k$}
\begin{lstlisting}
self._pi = N_k / N  # (K,)
\end{lstlisting}
Implementa la~\eqref{eq:pi_update}. La frazione di campioni (soft)
assegnata a ogni componente diventa il nuovo peso.

\subsubsection{Riga 5: aggiornamento medie $\mu_k$}
\begin{lstlisting}
self._mu = gamma.T @ samples / N_k[:, None]  # (K, d)
\end{lstlisting}

Questa riga è elegante. \texttt{gamma.T @ samples} è un prodotto
matriciale:
\[
  (\gamma^\top X)_{k,i} = \sum_{j=1}^{N} \gamma_{jk} \cdot x_{j,i}
\]
Shape: $(K, N) \times (N, d) = (K, d)$.

Dividendo per $N_k$ (con broadcasting da $(K,)$ a $(K,1)$ tramite
\texttt{[:, None]}), otteniamo la media pesata~\eqref{eq:mu_update}.

\subsubsection{Righe 6--8: aggiornamento varianze $\sigma_k^2$}
\begin{lstlisting}
diff = samples[:, None, :] - self._mu[None, :, :]  # (N, K, d)
self._sigma = (np.sum(gamma[:, :, None] * diff**2, axis=0)
               / N_k[:, None])                     # (K, d)
\end{lstlisting}

Passo per passo:
\begin{enumerate}[nosep]
  \item \texttt{diff}: differenze $x_{j,i} - \mu_{k,i}^{\text{new}}$
        (stesso broadcasting della Sezione~\ref{sec:logweighted}).
        Shape: $(N, K, d)$.
  \item \texttt{diff**2}: scarti al quadrato. Shape: $(N, K, d)$.
  \item \texttt{gamma[:, :, None]}: aggiunge un asse per le features,
        portando $\gamma$ da $(N, K)$ a $(N, K, 1)$ per il broadcasting.
  \item \texttt{gamma[:,:,None] * diff**2}: pesa ogni scarto quadrato
        con la responsabilità. Shape: $(N, K, d)$.
  \item \texttt{np.sum(..., axis=0)}: somma sui campioni. Shape: $(K, d)$.
  \item \texttt{/ N\_k[:, None]}: normalizza. Shape: $(K, d)$.
\end{enumerate}

Il risultato è esattamente la~\eqref{eq:sigma_diag}.

\subsubsection{Riga 9: floor sulle varianze}
\begin{lstlisting}
self._sigma = np.maximum(self._sigma, 1e-6)
\end{lstlisting}

\begin{warnbox}[Perche il floor sulle varianze]
Se una componente gaussiana collassa su un singolo punto ($\sigma_{k,i}^2
\to 0$), la densità $\mathcal{N}$ diverge a $+\infty$ e la
log-likelihood esplode. Questo fenomeno si chiama \textbf{singolarità}
del GMM. Il floor a $10^{-6}$ previene questo collasso garantendo che
ogni gaussiana mantenga una larghezza minima.
\end{warnbox}

\subsection{Convergenza}

\begin{lstlisting}[caption={Controllo di convergenza}]
ll = np.sum(log_denominator)
self._ll_history.append(ll)
if np.abs(ll - prev_ll) < self.epsilon:
    break
prev_ll = ll
\end{lstlisting}

Il criterio di convergenza è semplice: se la variazione assoluta
della log-likelihood tra due iterazioni consecutive è inferiore a
$\varepsilon = 10^{-8}$, l'algoritmo si ferma.

\begin{infobox}[Teorema di convergenza monotona --- GDL8 slide 15]
Un risultato fondamentale dell'EM è che la log-likelihood
\textbf{non decresce mai}:
\[
  \ell(\theta^{(t+1)} \mid X) \ge \ell(\theta^{(t)} \mid X)
\]
Graficamente (\refslide{GDL8}{15}), la curva della log-likelihood
è monotonamente crescente. Questo non garantisce il massimo globale,
ma garantisce che ogni iterazione migliora (o mantiene) la soluzione.

Nel codice, \texttt{self.\_ll\_history} registra la storia della
log-likelihood. Verificare che sia monotonamente crescente è un
ottimo test di correttezza (vedi Sezione~\ref{sec:checklist}).
\end{infobox}

\begin{exbox}[Nota su log\_denominator]
Nel ciclo EM, la variabile \texttt{log\_denominator} calcolata
nell'E-step è esattamente
$\log\sum_k \pi_k\mathcal{N}(x_j|\mu_k,\Sigma_k)$ per ogni $j$.
Quindi \texttt{np.sum(log\_denominator)} è la incomplete
log-likelihood $\ell(\theta^{(t)}|X)$. Non serve chiamare
\texttt{log\_likelihood()} separatamente: il valore è già
disponibile come sottoprodotto dell'E-step.
\end{exbox}

% ============================================================
% SECTION 9: BIC
% ============================================================
\section{BIC (Bayesian Information Criterion)}\label{sec:bic}

\subsection{Formula}

Il BIC è un criterio per la selezione del modello che bilancia la
bontà di adattamento (log-likelihood) con la complessità del modello
(numero di parametri):
\begin{equation}\label{eq:bic}
  \mathrm{BIC} = \ell(\hat\theta \mid X) -
  \frac{p}{2}\log N
\end{equation}
dove:
\begin{itemize}[nosep]
  \item $\ell(\hat\theta \mid X)$: log-likelihood massimizzata
        (incomplete).
  \item $p$: numero di parametri liberi del modello.
  \item $N$: numero di campioni.
\end{itemize}

Si sceglie il modello con il \textbf{BIC più alto} (massima likelihood
penalizzata).

\begin{warnbox}[Attenzione alla convenzione del segno]
Esistono due convenzioni per il BIC. Quella usata nel testo del midterm
(e nel nostro codice) è con il segno meno davanti alla penalità, quindi
si \textbf{massimizza}. In alcune fonti il BIC è definito con segno
opposto e si minimizza. Controlla sempre quale convenzione è usata!
\end{warnbox}

\subsection{Conteggio dei parametri}

Per un GMM con $K$ componenti, $d$ features, e covarianza diagonale:

\begin{center}
\renewcommand{\arraystretch}{1.3}
\begin{tabular}{lcc}
\toprule
\textbf{Parametro} & \textbf{Conteggio per componente} &
\textbf{Totale} \\
\midrule
$\pi_k$ (pesi) & --- & $K - 1$\footnotemark \\
$\mu_k$ (medie) & $d$ per componente & $K \cdot d$ \\
$\sigma_k^2$ (varianze diag.) & $d$ per componente & $K \cdot d$ \\
\midrule
\textbf{Totale} $p$ & & $(K-1) + 2Kd$ \\
\bottomrule
\end{tabular}
\end{center}
\footnotetext{I pesi hanno il vincolo $\sum_k \pi_k = 1$, quindi solo
$K-1$ sono liberi.}

Nel nostro caso con $d=5$:

\begin{center}
\begin{tabular}{cccc}
\toprule
$K$ & $K-1$ & $2Kd$ & $p = (K-1)+2Kd$ \\
\midrule
1 & 0 & 10 & 10 \\
2 & 1 & 20 & 21 \\
3 & 2 & 30 & 32 \\
4 & 3 & 40 & 43 \\
5 & 4 & 50 & 54 \\
6 & 5 & 60 & 65 \\
\bottomrule
\end{tabular}
\end{center}

\subsection{Il codice}

\begin{lstlisting}[caption={Calcolo del BIC}]
def bic(self, samples: np.ndarray) -> np.ndarray:
    ll = self.log_likelihood(samples)
    n = samples.shape[0]
    d = samples.shape[1]
    n_params = (self.n_categories - 1) + 2 * self.n_categories * d
    return ll - (n_params / 2) * np.log(n)
\end{lstlisting}

\begin{itemize}[nosep]
  \item \textbf{Riga 2}: calcola $\ell(\hat\theta|X)$ con la funzione
        già analizzata.
  \item \textbf{Riga 5}: conta i parametri: $(K-1) + 2Kd$.
  \item \textbf{Riga 6}: applica la formula~\eqref{eq:bic}.
\end{itemize}

\subsection{Risultati del training e interpretazione}

\begin{lstlisting}[language={},caption={Output del training}]
k=1  BIC=-8320.1771  logP(X|theta)=-8286.7541
k=2  BIC=-8052.0416  logP(X|theta)=-7981.8532
k=3  BIC=-7919.3844  logP(X|theta)=-7812.4306
k=4  BIC=-7869.8688  logP(X|theta)=-7726.1496
k=5  BIC=-7899.1658  logP(X|theta)=-7718.6813
k=6  BIC=-7916.2819  logP(X|theta)=-7699.0320
Best model: k=4
\end{lstlisting}

\begin{exbox}[Analisi dei risultati]
\textbf{Osservazioni}:
\begin{enumerate}[nosep]
  \item La log-likelihood $\ell(\hat\theta|X)$ cresce monotonamente
        con $K$ (da $-8286.8$ a $-7699.0$). Questo è atteso: più
        componenti = più flessibilità = migliore fit.
  \item Il BIC cresce da $K=1$ a $K=4$ e poi \textbf{cala} per $K=5$
        e $K=6$. La penalità sulla complessità supera il guadagno
        in likelihood.
  \item Il massimo BIC si ha per $K=4$ ($\mathrm{BIC}=-7869.87$),
        quindi il modello selezionato ha \textbf{4 componenti}.
\end{enumerate}

\textbf{Interpretazione}: il dataset contiene probabilmente 4 cluster
naturali. Con $K=5$ e $K=6$ si stanno modellando rumore o
sotto-strutture irrilevanti, e la penalità BIC li punisce.

\textbf{Penalità BIC per $K=4$ vs $K=5$}:
\begin{itemize}[nosep]
  \item $K=4$: penalità $= \frac{43}{2}\log(800) = 21.5 \times 6.685
        = 143.72$
  \item $K=5$: penalità $= \frac{54}{2}\log(800) = 27.0 \times 6.685
        = 180.49$
  \item La differenza di penalità ($\approx 36.8$) è maggiore del
        guadagno in log-likelihood ($-7718.7 - (-7726.1) = 7.5$),
        quindi $K=5$ non è giustificato.
\end{itemize}
\end{exbox}

% ============================================================
% SECTION 10: CHECKLIST
% ============================================================
\section{Checklist di verifica}\label{sec:checklist}

Questa checklist permette di verificare che l'implementazione sia
corretta e che il modello funzioni come atteso.

\subsection{Proprietà da verificare}

\begin{enumerate}[label=\textbf{\arabic*.}, leftmargin=2em]

\item \textbf{La log-likelihood cresce monotonamente}

Dopo il training, verifica che \texttt{self.\_ll\_history} sia una
sequenza non decrescente:
\begin{lstlisting}[numbers=none]
diffs = np.diff(model._ll_history)
assert np.all(diffs >= -1e-6), "LL non monotona!"
\end{lstlisting}
Piccole violazioni ($\sim 10^{-10}$) sono accettabili per errori di
arrotondamento. Violazioni significative indicano un bug.

\item \textbf{Le righe di gamma sommano a 1}

Dopo l'E-step, ogni $\gamma_{j,:}$ deve sommare a~1 (è una
distribuzione di probabilità sulle componenti):
\begin{lstlisting}[numbers=none]
row_sums = gamma.sum(axis=1)
assert np.allclose(row_sums, 1.0), "Gamma rows != 1"
\end{lstlisting}

\item \textbf{I pesi pi sommano a 1}

I mixing coefficients devono formare una distribuzione:
\begin{lstlisting}[numbers=none]
assert np.isclose(model._pi.sum(), 1.0), "Pi sum != 1"
\end{lstlisting}

\item \textbf{Le varianze sono sempre positive}

Grazie al floor \texttt{np.maximum(self.\_sigma, 1e-6)}, le varianze
non possono diventare negative o zero:
\begin{lstlisting}[numbers=none]
assert np.all(model._sigma > 0), "Varianze non positive!"
\end{lstlisting}

\item \textbf{Il BIC seleziona correttamente $K$}

Il modello con il BIC più alto deve essere quello ottimale. Nel nostro
caso, $K=4$ deve avere il BIC massimo tra tutti i modelli provati.
\begin{lstlisting}[numbers=none]
bics = [models[k].bic(train_data) for k in range(1,7)]
best_k = np.argmax(bics) + 1
assert best_k == 4, f"BIC seleziona K={best_k}, atteso K=4"
\end{lstlisting}

\item \textbf{Assenza di NaN e Inf}

Nessun parametro deve contenere valori non finiti:
\begin{lstlisting}[numbers=none]
assert np.all(np.isfinite(model._mu)), "NaN/Inf in mu"
assert np.all(np.isfinite(model._sigma)), "NaN/Inf in sigma"
assert np.all(np.isfinite(model._pi)), "NaN/Inf in pi"
assert np.isfinite(model.log_likelihood(train_data)), "LL non finita"
\end{lstlisting}

\item \textbf{Le sezioni 3 e 4 del notebook non sono state alterate}

Il testo del midterm specifica che alcune sezioni del notebook non devono
essere modificate. Verificare che le celle delle sezioni~3 e~4 siano
rimaste invariate rispetto all'originale.
\begin{lstlisting}[numbers=none]
# Verifica manuale: confronta le celle del notebook
# con la versione originale scaricata
\end{lstlisting}

\end{enumerate}

\subsection{Tabella riassuntiva}

\begin{center}
\renewcommand{\arraystretch}{1.3}
\begin{tabular}{clcc}
\toprule
\textbf{\#} & \textbf{Proprietà} & \textbf{Dove verificare} &
\textbf{Tolleranza} \\
\midrule
1 & Log-likelihood monotona &
\texttt{\_ll\_history} & $10^{-6}$ \\
2 & $\sum_k \gamma_{jk} = 1$ &
dopo E-step & $10^{-10}$ \\
3 & $\sum_k \pi_k = 1$ &
\texttt{\_pi} & $10^{-10}$ \\
4 & $\sigma_{k,i}^2 > 0$ &
\texttt{\_sigma} & $> 10^{-6}$ \\
5 & BIC $\to K=4$ &
output training & esatto \\
6 & No NaN/Inf &
tutti i parametri & esatto \\
7 & Sez.\ 3--4 invariate &
notebook & esatto \\
\bottomrule
\end{tabular}
\end{center}

\begin{tipbox}[Consiglio pratico per il debug]
Se qualcosa non funziona, controlla in quest'ordine:
\begin{enumerate}[nosep]
  \item Stampa le shape di tutti gli array intermedi e confrontale
        con la tabella della Sezione~\ref{sec:logweighted}.
  \item Verifica che la log-likelihood cresca (punto~1).
  \item Controlla $\gamma$ (punto~2): se non somma a~1, il problema
        è nel log-sum-exp.
  \item Controlla le varianze (punto~4): se diventano troppo piccole,
        il modello sta collassando.
\end{enumerate}
\end{tipbox}

% ============================================================
\newpage
\section*{Riepilogo dei riferimenti alle slide}

\begin{center}
\renewcommand{\arraystretch}{1.3}
\begin{tabular}{ll}
\toprule
\textbf{Riferimento} & \textbf{Contenuto} \\
\midrule
GDL7 slide 9 & Tabella apprendimento: struttura vs dati \\
GDL8 slide 3 & Stessa tabella con EM evidenziato \\
GDL8 slide 13 & Incomplete data e variabili latenti \\
GDL8 slide 14 & Schema generale EM (E-step / M-step) \\
GDL8 slide 15 & Convergenza monotona di EM (grafico) \\
GDL8 slide 18 & Formula gaussiana univariata \\
GDL8 slides 19--20 & Definizione GMM \\
GDL8 slide 21 & Incomplete likelihood \\
GDL8 slide 24 & Complete likelihood con variabili indicatrici \\
GDL8 slide 25 & E-step: formula delle responsabilità \\
GDL8 slide 27 & M-step: formule di aggiornamento \\
\bottomrule
\end{tabular}
\end{center}

\end{document}
